El grafo asociado a un problema de transporte con cinco plantas (orígenes) y cinco almacenes (destinos) es un grafo dirigido. Si el mismo grafo fuera no dirigido:
\begin{itemize}
    \item Indica un ciclo euleriano en dicho grafo no dirigido. \textit{Un grafo conexo no dirigido contiene un ciclo euleriano si y solo si todos sus vértices tienen grado par. En el grafo no dirigido todos los nodos tienen grado 5; luego   no contiene un ciclo euleriano.}
    \item Indica un ciclo hamiltoniano en dicho grafo no dirigido.  \textit{Existen  muchos ciclos hamiltonianos. Si se llaman $O_i$ los nodos origen y $D_j$ los nodos destino, un posible ciclo hamiltoniano sería: $O_1 \rightarrow D_1 \rightarrow O_2 \rightarrow D_2 \rightarrow O_3 \rightarrow D_3 \rightarrow O_4 \rightarrow D_4 \rightarrow O_5 \rightarrow D_5 \rightarrow O_1$}
    \item ¿Existen ciclos euleriano para cualquier número de orígenes y destinos? \textit{Si el número de orígenes y de destinos es par, existirán ciclos eulerianos.} 
    \item ¿Existen ciclos hamiltonianos para cualquier número de orígenes y destinos? \textit{Si el número de orígenes y destinos es el mismo, existirán ciclos hamiltoniano  (y al menos hay dos orígenes y dos destinos).}


\end{itemize}

